% ============================================================================
% Fichier  : front/01_avant_propos.tex
% Titre    : Avant-propos
% Auteur   : MINKA MI NGUIDJOÏ Thierry Emmanuel
% Version  : 1.0 -- novembre 2026
% ============================================================================

\chapter*{Avant-propos}
\addcontentsline{toc}{chapter}{Avant-propos}
\markboth{Avant-propos}{Avant-propos}

\epigraph{« On ne protège bien que ce que l'on connaît,\\
et l'on ne connaît bien que ce que l'on a mesuré. »}{Principe du cours}

Ce cours s'adresse aux élèves ingénieurs de troisième année du Département
de Génie Informatique de l'École Nationale Supérieure Polytechnique de
Yaoundé. Il est leur premier cours complet de sécurité informatique~: il
suppose des bases en systèmes, en réseaux et en programmation, mais aucune
connaissance préalable de la sécurité. Il ouvre la voie aux UE d'audit
(INF3226), d'investigation numérique (SEC4031 et SEC4052), d'analyse des
risques et de gestion de crise (GES4062) et de hacking éthique (SEC4062).

Il repose sur une conviction~: la sécurité n'est pas un produit, mais une
discipline. On ne l'achète pas avec un pare-feu ou un antivirus~; on la
construit en sachant ce que l'on protège, contre qui, avec quelles
mesures, et en vérifiant qu'elles fonctionnent. Les attaques réelles
réussissent rarement par une technique nouvelle~; elles réussissent parce
qu'une mesure connue manquait, ou qu'on la croyait en place sans l'avoir
vérifiée.

\bigskip

\textbf{Ce que ce cours n'est pas.} Il n'est pas une préparation à une
certification, ni un manuel d'outils, ni un cours d'attaque. Les
techniques offensives n'y sont présentées que pour comprendre ce dont on
se protège, et toujours dans le cadre légal camerounais.

\bigskip

\textbf{Ce que ce cours est.} Un traité de vingt-six chapitres, en six
parties précédées d'un prologue. La Partie~I pose les fondements~:
propriétés de sécurité, menaces et attaques, cryptographie, grandes
attaques de l'histoire. La Partie~II traite de la gouvernance, des
risques, des actifs et du facteur humain. La Partie~III présente les
normes, le droit camerounais et africain, et l'éthique. La Partie~IV
couvre la sécurité technique~: modèles formels, architecture, réseaux,
identités, systèmes, applications, nuage, évaluation et opérations. La
Partie~V aborde la continuité, l'investigation numérique et les
technologies émergentes. La Partie~VI apprend à mesurer la sécurité, et à
dire ce que l'on ne sait pas.

\bigskip

\textbf{Un fil rouge.} Le cours suit une entreprise fictive, LiZbiZ,
société de logistique aux quatre centres de données, attaquée par un
rançongiciel une nuit de mars (Opération GAROUA). Chaque chapitre revient
sur cette nuit pour montrer ce qui a manqué et ce qui aurait changé. La
BANQUE SANAGA et l'ANPP, agence publique de paiement, complètent le
tableau. Les noms, les dates et les montants sont fictifs~; les
mécanismes d'attaque et de défense sont exacts, et les textes de loi sont
cités d'après leur version officielle.

\bigskip

\textbf{Comment travailler.} Lisez chaque chapitre avec sa fiche
d'auto-évaluation~; faites les exercices avant de lire leurs corrigés~;
pratiquez dans les labs de NDAP SUKULU. Prenez l'habitude que le cours
cherche à transmettre~: séparer ce que l'on a constaté de ce que l'on en
pense, et ne jamais présenter comme sûr ce que l'on n'a pas vérifié.

\bigskip

\begin{flushright}
    \IfFileExists{../../Logos/signature_avant_propos.png}{\includegraphics[height=1.6cm]{../../Logos/signature_avant_propos.png}}{\textit{MINKA MI NGUIDJOÏ Thierry Emmanuel}}\\
    \textit{Yaoundé, novembre 2026}
\end{flushright}
